\chapter{Personae et parcours utilisateurs}\label{annexe:persona}



\subsection{Présentation des personae}

Les personae représentent des profils types d’utilisateurs du projet \emph{5 Secondes Chrono}.
Ils permettent d’orienter les choix fonctionnels, l’ergonomie et les priorités du produit
en s’appuyant sur des besoins réalistes.

Chaque persona comprend :
\begin{itemize}
    \item un nom (fictif) et une tranche d’âge ;
    \item un contexte d’utilisation ;
    \item des motivations et objectifs ;
    \item des frustrations (points de douleur) ;
    \item une citation synthétisant son état d’esprit.
\end{itemize}

\subsection{Personae retenus pour \emph{5 Secondes Chrono}}

\begin{table}[h!]
\centering
\renewcommand{\arraystretch}{1.2}
\small
\begin{tabular}{|p{2.6cm}|p{3.2cm}|p{4.9cm}|p{4.9cm}|}
\hline
\rowcolor{etnaBlue!15}
\textbf{Persona} & \textbf{Profil / Contexte} & \textbf{Motivations / Objectifs} & \textbf{Frustrations} \\ \hline

\textbf{Nicolas (32 ans)} &
Négociateur en immobilier commercial. Utilise l’app le matin ou entre deux RDV. &
Monter en compétences rapidement ; réviser des notions clés ; progresser sans y passer 1h. &
Manque de temps ; contenus trop denses ; outils non adaptés au mobile ; peur de ``perdre son niveau''. \\ \hline

\textbf{Sarah (27 ans)} &
Jeune professionnelle en prise de poste. Utilise l’app le soir en micro-sessions. &
Se sentir à l’aise sur les bases juridiques ; apprendre de manière structurée ; avoir un feedback motivant. &
Stress d’erreur ; manque d’accompagnement ; difficulté à savoir quoi réviser ; besoin d’encouragements. \\ \hline

\textbf{Karim (41 ans)} &
Indépendant, agenda chargé. Utilise l’app de façon irrégulière. &
Réviser ponctuellement ; accéder rapidement à une explication (vidéo) ; gagner en confiance. &
Oublie de revenir ; n’aime pas les interfaces complexes ; abandon si la valeur n’est pas immédiate. \\ \hline

\textbf{Claire (35 ans)} &
Administratrice / formateur côté \emph{Le Carré Pro}. Gère contenus et suivi. &
Ajouter/mettre à jour les questions/vidéos ; suivre la progression ; s’assurer de la qualité pédagogique. &
Processus de gestion fastidieux ; manque de visibilité sur l’engagement ; crainte d’erreurs de contenu. \\ \hline

\end{tabular}
\caption{Personae principaux du projet \emph{5 Secondes Chrono}}
\end{table}

\subsection{Parcours utilisateurs}

Les parcours utilisateurs décrivent les étapes suivies par un persona pour atteindre un objectif.
Ils permettent d’identifier les points de friction, les besoins d’assistance et les priorités
d’interface.

Le projet \emph{5 Secondes Chrono} suit un parcours apprenant type :
\emph{Accueil $\rightarrow$ Quiz du jour $\rightarrow$ Résultats $\rightarrow$ Progression $\rightarrow$ Recommandations}.
Ce flux correspond au parcours macro défini pour l’application. \par
\noindent\textit{(Référence : séquence fonctionnelle globale du projet)} \textbf{---} \emph{Accueil / Quiz / Résultats / Progression / Reco}. \hfill

\subsection{Parcours utilisateur 1 — Réaliser le quiz quotidien (Nicolas)}

\begin{table}[h!]
\centering
\renewcommand{\arraystretch}{1.15}
\small
\begin{tabular}{|p{3.2cm}|p{10.4cm}|}
\hline
\rowcolor{black!10}
\textbf{Étape} & \textbf{Action / Résultat attendu} \\ \hline
Point d’entrée & Nicolas ouvre l’application depuis son téléphone (session courte). \\ \hline
Quiz du jour & Il lance le quiz quotidien et répond aux questions avec le timer. \\ \hline
Résultats & Il consulte son score et les explications associées (dont vidéo si disponible). \\ \hline
Progression & Il visualise son niveau / sa progression (thèmes maîtrisés vs à revoir). \\ \hline
Recommandations & Il reçoit une recommandation simple (notion à revoir / contenu conseillé). \\ \hline
Fin & Il quitte l’application en ayant le sentiment d’avoir progressé en moins de 2 minutes. \\ \hline
\end{tabular}
\caption{Parcours utilisateur — Quiz quotidien}
\end{table}

\subsection{Parcours utilisateur 2 — Se remettre à niveau grâce aux recommandations (Sarah)}

\begin{table}[h!]
\centering
\renewcommand{\arraystretch}{1.15}
\small
\begin{tabular}{|p{3.2cm}|p{10.4cm}|}
\hline
\rowcolor{black!10}
\textbf{Étape} & \textbf{Action / Résultat attendu} \\ \hline
Point d’entrée & Sarah ouvre l’app après une journée de travail, objectif : réviser ``utile''. \\ \hline
Progression & Elle consulte ses statistiques et repère une thématique en difficulté. \\ \hline
Recommandations & L’app suggère un quiz ciblé / une vidéo liée à ses erreurs fréquentes. \\ \hline
Contenu & Elle regarde une vidéo courte ou lit une explication avant de relancer un quiz. \\ \hline
Résultats & Elle vérifie l’amélioration de son score sur le thème concerné. \\ \hline
Fin & Elle ressort rassurée, avec une action claire pour la prochaine session. \\ \hline
\end{tabular}
\caption{Parcours utilisateur — Remédiation par recommandations}
\end{table}

\subsection{Parcours utilisateur 3 — Usage irrégulier mais efficace (Karim)}

\begin{table}[h!]
\centering
\renewcommand{\arraystretch}{1.15}
\small
\begin{tabular}{|p{3.2cm}|p{10.4cm}|}
\hline
\rowcolor{black!10}
\textbf{Étape} & \textbf{Action / Résultat attendu} \\ \hline
Point d’entrée & Karim relance l’app après plusieurs jours sans usage. \\ \hline
Accueil & L’app propose un accès direct au quiz et un rappel rapide de sa progression. \\ \hline
Quiz du jour & Il réalise une session courte, sans configuration complexe. \\ \hline
Résultats & Il comprend immédiatement ce qu’il a raté et pourquoi (explication claire). \\ \hline
Fin & Il obtient une valeur immédiate et est incité à revenir (gamification légère). \\ \hline
\end{tabular}
\caption{Parcours utilisateur — Reprise rapide}
\end{table}

\subsection{Parcours utilisateur 4 — Gestion des contenus (Claire, administratrice)}

\begin{table}[h!]
\centering
\renewcommand{\arraystretch}{1.15}
\small
\begin{tabular}{|p{3.2cm}|p{10.4cm}|}
\hline
\rowcolor{black!10}
\textbf{Étape} & \textbf{Action / Résultat attendu} \\ \hline
Point d’entrée & Claire se connecte à l’espace administrateur (web). \\ \hline
Gestion contenus & Elle ajoute/édite une question et associe une vidéo à une thématique. \\ \hline
Contrôle qualité & Elle vérifie la cohérence et la formulation (validation interne). \\ \hline
Publication & Les contenus deviennent disponibles selon les règles d’accès (abonnement). \\ \hline
Suivi & Elle consulte des statistiques (engagement, difficultés fréquentes). \\ \hline
Fin & Elle ajuste le contenu pédagogique selon les retours et les données observées. \\ \hline
\end{tabular}
\caption{Parcours utilisateur — Administration des contenus}
\end{table}

